% Relatorio do Projeto AM 2026-2 (CIn/UFPE).
% Compile neste diretorio com: latexmk -pdf relatorio.tex
% Todos os graficos sao desenhados com pgfplots: os da Questao 1 e as curvas de aprendizagem
% a partir de ../../results, e os boxplots a partir de dados/boxplots.csv (gerado por gerar_boxplots.py).
\documentclass[11pt,a4paper]{article}

\usepackage[T1]{fontenc}
\usepackage[utf8]{inputenc}
% English is loaded only for the IEEE reference strings; Portuguese is the main language.
\usepackage[american,brazilian]{babel}
\usepackage{lmodern}
\usepackage{microtype}
\usepackage[a4paper,margin=2.5cm]{geometry}
\usepackage{amsmath,amssymb,bm}
\usepackage{booktabs,multirow,array,tabularx}
\usepackage{siunitx}
\sisetup{
  output-decimal-marker = {,},
  group-separator = {.},
  group-minimum-digits = 4,
  group-digits = integer,
  mode = match,
  propagate-math-font = true,
  reset-math-version = false,
  reset-text-series = false,
  text-series-to-math = true,
}
\usepackage{graphicx}
\usepackage{caption,subcaption}
\captionsetup{font=small,labelfont=bf}
\usepackage{xcolor}
\usepackage{tikz}
\usetikzlibrary{positioning,arrows.meta,calc,decorations.pathreplacing}
\usepackage{pgfplots}
\pgfplotsset{compat=1.18}
\usepgfplotslibrary{groupplots,fillbetween,statistics}
% Keeps only the table rows whose column #1 equals the string #2.
\pgfplotsset{discard if not/.style 2 args={x filter/.append code={%
  \ifnum\pdfstrcmp{\thisrow{#1}}{#2}=0 \else\def\pgfmathresult{inf}\fi}}}
\usepackage{float}
\usepackage{enumitem}
\setlist{itemsep=2pt,topsep=4pt}
\usepackage{csquotes}
% Numbered citations and reference list in the IEEE style (biblatex-ieee).
\usepackage[backend=biber,style=ieee,maxbibnames=6,minbibnames=1]{biblatex}
\addbibresource{referencias.bib}
\usepackage[hidelinks]{hyperref}
\usepackage{tocloft}
\renewcommand{\cfttoctitlefont}{\hfill\bfseries\MakeUppercase}
\renewcommand{\cftaftertoctitle}{\hfill\mbox{}}
% The contents page is pre-textual: no page number (tocloft would force 'plain').
\tocloftpagestyle{empty}

% Wider gaps between table columns (default: 6pt).
\setlength{\tabcolsep}{8pt}
% Ragged-right X column, so narrow cells are not stretched by justification.
\newcolumntype{L}{>{\raggedright\arraybackslash}X}

% Media com meia-amplitude do IC 95%.
\newcommand{\ci}[2]{\num{#1}\,{\footnotesize$\pm$\,\num{#2}}}
% Melhor valor de uma coluna.
\newcommand{\best}[1]{\textbf{#1}}
% p-valor de Nemenyi nao significativo (alpha = 0,05).
\newcommand{\ns}[1]{\textbf{\num{#1}}}
% F-measure de treino / teste na curva de aprendizagem.
\newcommand{\lc}[2]{\num{#1}\,/\,\num{#2}}
% Passage still to be written by the team (red box).
\newcommand{\todoequipe}[1]{%
  \par\medskip\noindent\fcolorbox{red!70!black}{red!4}{\parbox{\dimexpr\linewidth-2\fboxsep-2\fboxrule\relax}{\small\textbf{A completar:} #1}}\par\medskip}

% Learning-curve panel read from the summary CSV: #1 dataset, #2 classifier, #3 panel title.
% Lines are the mean of the 10 splits; bands are +-1 standard deviation, capped at 1.
\newcommand{\lcfile}{../../results/learning_curves_summary.csv}
\newcommand{\lcband}[4]{% #1 dataset, #2 classifier, #3 split (train/test), #4 colour
  \addplot[name path=up#3, draw=none, forget plot, discard if not={dataset}{#1}, discard if not={classifier}{#2}]
    table[col sep=comma, x expr=100*\thisrow{train_size},
          y expr={min(\thisrow{#3_f1_mean}+\thisrow{#3_f1_std},1)}] {\lcfile};
  \addplot[name path=lo#3, draw=none, forget plot, discard if not={dataset}{#1}, discard if not={classifier}{#2}]
    table[col sep=comma, x expr=100*\thisrow{train_size},
          y expr={\thisrow{#3_f1_mean}-\thisrow{#3_f1_std}}] {\lcfile};
  \addplot[#4, fill opacity=0.18, draw=none, forget plot] fill between[of=up#3 and lo#3];}
\newcommand{\lcpanel}[3]{%
  \nextgroupplot[title={#3}]
  \lcband{#1}{#2}{train}{blue!70!black}
  \lcband{#1}{#2}{test}{orange!90!black}
  \addplot[blue!70!black, thick, mark=*, mark size=1.1pt, discard if not={dataset}{#1}, discard if not={classifier}{#2}]
    table[col sep=comma, x expr=100*\thisrow{train_size}, y=train_f1_mean] {\lcfile};
  \addplot[orange!90!black, thick, mark=square*, mark size=1.1pt, discard if not={dataset}{#1}, discard if not={classifier}{#2}]
    table[col sep=comma, x expr=100*\thisrow{train_size}, y=test_f1_mean] {\lcfile};}
% Full learning-curve figure body: #1 dataset, #2 ymin, #3 y ticks.
\newcommand{\lcfigure}[3]{%
\begin{tikzpicture}
\begin{groupplot}[
  group style={group size=2 by 3, horizontal sep=1.5cm, vertical sep=1.9cm, ylabels at=edge left},
  scale only axis, width=6.3cm, height=3.9cm,
  xmin=2, xmax=98, xtick={5,25,50,75,95},
  ymin=#2, ymax=1.02, ytick={#3},
  xlabel={Treino (\% da base)}, ylabel={F-measure macro},
  grid=major, grid style={dashed, gray!30},
  title style={font=\small\bfseries, yshift=-0.6ex},
  tick label style={font=\footnotesize}, label style={font=\footnotesize},
  yticklabel style={/pgf/number format/.cd, fixed, fixed zerofill, precision=1, use comma},
]
\lcpanel{#1}{GaussianBayes}{Bayesiano gaussiano}
\lcpanel{#1}{BayesianKNN}{Bayesiano $k$-NN}
\lcpanel{#1}{ParzenWindow}{Janela de Parzen}
\lcpanel{#1}{LogisticRegression}{Regressão logística}
\lcpanel{#1}{MajorityVoting}{Voto majoritário}
\nextgroupplot[hide axis, xmin=0, xmax=1, ymin=0, ymax=1,
  legend style={at={(0.5,0.5)}, anchor=center, draw=gray!60, font=\small, row sep=2pt},
  legend cell align=left]
\addlegendimage{blue!70!black, thick, mark=*, mark size=1.3pt}\addlegendentry{Treino (média)}
\addlegendimage{area legend, draw=none, fill=blue!70!black, fill opacity=0.18}\addlegendentry{Treino $\pm$ 1 desvio-padrão}
\addlegendimage{orange!90!black, thick, mark=square*, mark size=1.3pt}\addlegendentry{Teste (média)}
\addlegendimage{area legend, draw=none, fill=orange!90!black, fill opacity=0.18}\addlegendentry{Teste $\pm$ 1 desvio-padrão}
\end{groupplot}
\end{tikzpicture}}

% Boxplots of the 300 test folds (Figure 5), from statistics computed by gerar_boxplots.py.
% Rows are ordered by metric, dataset and classifier: panel k uses rows 5k to 5k+4.
\pgfplotstableread[col sep=comma]{dados/boxplots.csv}\bptable
\definecolor{cGB}{HTML}{1F77B4}\definecolor{cKNN}{HTML}{FF7F0E}\definecolor{cPZ}{HTML}{2CA02C}
\definecolor{cRL}{HTML}{D62728}\definecolor{cVM}{HTML}{9467BD}
\newcommand{\bpget}[3]{\pgfplotstablegetelem{#1}{#2}\of\bptable\let#3\pgfplotsretval}
% One box from table row #1 with fill colour #2. The \edef writes the numbers into the
% \addplot options, so each box keeps its own values.
\newcommand{\bpbox}[2]{%
  \bpget{#1}{lower_whisker}\bpLW\bpget{#1}{q1}\bpQA\bpget{#1}{median}\bpMed
  \bpget{#1}{q3}\bpQC\bpget{#1}{upper_whisker}\bpUW\bpget{#1}{mean}\bpMean
  \edef\bpcmd{\noexpand\addplot[draw=black, fill=#2!65, boxplot prepared={lower whisker=\bpLW,
    lower quartile=\bpQA, median=\bpMed, upper quartile=\bpQC, upper whisker=\bpUW,
    average=\bpMean}] coordinates {};}%
  \bpcmd}
\newcommand{\bppanel}[1]{% #1 = first table row of the panel
  \bpbox{#1}{cGB}\bpbox{\the\numexpr#1+1}{cKNN}\bpbox{\the\numexpr#1+2}{cPZ}%
  \bpbox{\the\numexpr#1+3}{cRL}\bpbox{\the\numexpr#1+4}{cVM}}

\newcommand{\Sil}{\mathrm{Sil}}
\newcommand{\ARI}{\mathrm{ARI}}

\hypersetup{
  pdftitle={Double k-means e classificadores bayesianos, logístico e por voto majoritário na base Spambase},
  pdfauthor={Bruna Isabel Santos Cruz, Everthon Manoel da Silva Inacio, Ana Carolina Soares de Melo},
}

\begin{document}
% Pre-textual pages (cover, abstract, contents) carry no page number but are counted.
\pagestyle{empty}

%----------------------------------------------------------------------
% CAPA 
%----------------------------------------------------------------------
\begin{titlepage}
  \centering
  \includegraphics[width=0.2\textwidth]{imagens/brasao.png}\par
  \vspace{0.6cm}
  UNIVERSIDADE FEDERAL DE PERNAMBUCO\par
  \vspace{0.2cm}
  CENTRO DE INFORMÁTICA\par
  \vspace{0.2cm}
  APRENDIZAGEM DE MÁQUINA -- 2026.2\par

  \vspace*{1.8cm}
  {\large BRUNA ISABEL SANTOS CRUZ\par
  \vspace{0.15cm}
  EVERTHON MANOEL DA SILVA INACIO\par
  \vspace{0.15cm}
  ANA CAROLINA SOARES DE MELO\par}

  \vspace*{2.4cm}
  {\Large\bfseries Double k-means e classificadores bayesianos, logístico e por voto majoritário na base
  Spambase\par}

  \vspace{1.6cm}
  \begin{flushright}
  \begin{minipage}{0.55\textwidth}
    \small
    Relatório do projeto da disciplina de Aprendizagem de Máquina do Centro de Informática da Universidade
    Federal de Pernambuco, elaborado por estudantes de mestrado do CIn/UFPE.
    \par\vspace{0.6cm}
    \textbf{Professor} Prof.\ Francisco de A.\ T.\ de Carvalho\par
  \end{minipage}
  \end{flushright}

  \vfill
  {\large Recife\par
  2026\par}
\end{titlepage}
% titlepage resets the counter to 1; the cover still counts as page 1.
\setcounter{page}{2}

%----------------------------------------------------------------------
% RESUMO
%----------------------------------------------------------------------
\clearpage
\begin{center}
  \bfseries\MakeUppercase{Resumo}
\end{center}
\vspace{0.5cm}

O presente relatório descreve os experimentos do projeto da disciplina de Aprendizagem de Máquina sobre a base Spambase ($\num{4601}$ e-mails, 57 atributos contínuos). 
Na Questão~1, o algoritmo double k-means foi executado 100 vezes para cada um dos 9 pares $(K,H)$ com $K\in\{2,3,4\}$ e $H\in\{1,\dots,K\}$; 
a melhor execução de cada par, segundo a função objetivo, foi comparada pela silhueta, que selecionou $(K^*,H^*)=(2,2)$ ($\Sil=\num{0.6596}$). 
A partição resultante isola 34 e-mails legítimos com valores extremos em 11 atributos ligados à HP e tem índice de Rand corrigido $\ARI=\num{-0.00495}$ 
em relação às classes a priori. Na Questão~2, cinco classificadores (bayesiano gaussiano, bayesiano $k$-NN, bayesiano com janela de Parzen, regressão 
logística e voto majoritário) foram avaliados nas duas versões da base (rótulo original e rótulo dado pelos $K^*$ grupos) por validação cruzada 
estratificada $30\times10$ com ajuste de hiperparâmetros por validação cruzada interna de 5 folds, intervalos de confiança, testes de Friedman e Nemenyi 
e curvas de aprendizagem. Na base original, voto majoritário ($F=\num{0.9238}$) e regressão logística ($F=\num{0.9230}$) não diferem significativamente 
entre si e superam os demais.

%----------------------------------------------------------------------
% SUMÁRIO
%----------------------------------------------------------------------
\clearpage
\tableofcontents

% Text pages: numbering continues from the pre-textual pages.
\clearpage
\pagestyle{plain}

%======================================================================
\section{Introdução}
%======================================================================

O projeto possui duas partes, sendo a primeira relativa à Questão~1, de natureza não supervisionada, na qual o double k-means~\cite{vichi2001double} 
é aplicado para agrupar ao mesmo tempo os objetos e as variáveis da Spambase, o número de grupos de objetos $K^*$ é escolhido pela 
silhueta~\cite{rousseeuw1987silhouettes} e a partição obtida é comparada com as classes a priori pelo índice de Rand corrigido~\cite{hubert1985comparing}. 
Já a segunta parte, que trata da Questão~2, de natureza supervisionada, onde cinco classificadores são avaliados e comparados em duas versões da base, 
uma com o rótulo original (spam/não spam) e outra com o rótulo dado pela partição de $K^*$ grupos da Questão~1.

Atendendo às observações finais do enunciado, o relatório dá destaque a três pontos: a organização dos experimentos que garante uma avaliação sem
vazamento de informação e uma comparação pareada entre os modelos (Seção~\ref{sec:organizacao}), a descrição dos dados (Seção~\ref{sec:dados}) e 
a obtenção dos hiperparâmetros (Seção~\ref{sec:hiper}). A Seção~\ref{sec:q1} trata da Questão~1, a Seção~\ref{sec:q2} da Questão~2, a
Seção~\ref{sec:discussao} discute limitações e a Seção~\ref{sec:conclusao} traz a conclusão do trabalho. Todo o código está no repositório do projeto, 
além disso, o Apêndice~\ref{app:reproducao} lista os comandos que geram cada resultado.

%======================================================================
\section{Descrição dos dados}\label{sec:dados}
%======================================================================

A base Spambase~\cite{hopkins1999spambase}, do repositório UCI, contém $n=\num{4601}$ e-mails, sem valores ausentes, descritos por uma variável 
resposta binária (1 = spam, 0 = não spam) e por $d=57$ atributos contínuos, que formam os três grupos da Tabela~\ref{tab:atributos}.

\begin{table}[H]
\centering
\caption{Grupos de atributos da Spambase.}\label{tab:atributos}
\small
\begin{tabularx}{\linewidth}{@{}l c L@{}}
\toprule
Grupo & Qtd. & Descrição \\
\midrule
\texttt{word\_freq\_*} & 48 & Percentual ($[0,100]$) de palavras do e-mail iguais a uma palavra dada
  (p.\,ex.\ \textit{make}, \textit{free}, \textit{money}, \textit{hp}, \textit{george}, \textit{650}). \\
\texttt{char\_freq\_*} & 6 & Percentual ($[0,100]$) de caracteres iguais a \texttt{;}, \texttt{(}, \texttt{[},
  \texttt{!}, \texttt{\$} e \texttt{\#}. \\
\texttt{capital\_run\_length\_*} & 3 & Comprimento médio (\textit{average}) e máximo (\textit{longest}) das
  sequências de letras maiúsculas e número total de maiúsculas (\textit{total}); valores $\ge 1$. \\
\bottomrule
\end{tabularx}
\end{table}

As características abaixo foram levantadas na base completa e orientaram decisões dos experimentos:

\begin{itemize}
\item \textbf{Classes moderadamente desbalanceadas:} $\num{2788}$ e-mails não spam (60,6\,\%) e $\num{1813}$ spam
  (39,4\,\%), razão pela qual todas as partições de treino e teste foram estratificadas.
\item \textbf{Escalas muito diferentes:} os desvios-padrão dos atributos vão de $\num{0.076}$ a $\num{606.3}$. As
  frequências ficam quase sempre abaixo de 10, enquanto \texttt{capital\_run\_length\_total} chega a
  $\num{15841}$ (mediana 95). Métodos baseados em distância ou com penalização dos coeficientes dependem da
  escala, o que motivou a padronização z-score na Questão~1 e dentro dos modelos $k$-NN e regressão logística na
  Questão~2.
\item \textbf{Esparsidade e assimetria:} 77,4\,\% das entradas da matriz são zero, já que 50 dos 57 atributos têm
  mais de metade dos valores iguais a zero e 19 têm mais de 90\,\%, e a assimetria amostral mediana dos atributos
  chega a 9,7. A hipótese de normalidade multivariada do classificador gaussiano é, portanto, pouco realista, e em amostras
  pequenas alguns atributos ficam constantes, o que torna as covariâncias singulares
  (Seção~\ref{sec:classificadores}).
\item \textbf{Linhas duplicadas:} há 391 linhas idênticas (atributos e rótulo) a alguma linha anterior. Como a base
  original foi mantida intacta, cópias idênticas podem cair uma no treino e outra no teste, o que favorece
  levemente os métodos de vizinhança (Seção~\ref{sec:discussao}).
\item \textbf{Origem dos e-mails legítimos:} segundo a documentação da base, os e-mails não spam vieram da
  correspondência de trabalho e pessoal dos autores na Hewlett-Packard, de modo que palavras como
  \textit{george} e o código de área \textit{650} indicam não spam~\cite{hopkins1999spambase}, fato que
  explica o grupo de variáveis encontrado na Questão~1.
\end{itemize}

A segunda versão da base, usada na Questão~2, mantém os mesmos 57 atributos e troca a variável resposta pela
partição de objetos com $K^*=2$ grupos obtida na Questão~1: $\num{4567}$ objetos (99,26\,\%) no grupo $P_1$ e
34 (0,74\,\%) no grupo $P_2$ (Seção~\ref{sec:particao}). Essa versão é extremamente desbalanceada: um classificador
que sempre prediz o grupo majoritário erra apenas 0,74\,\% dos exemplos.

%======================================================================
\section{Organização dos experimentos}\label{sec:organizacao}
%======================================================================

A Figura~\ref{fig:pipeline} resume o encadeamento dos experimentos: a Questão~1 usa apenas os atributos e produz os rótulos $y^*$ da segunda versão 
da base, enquanto e a Questão~2 avalia os classificadores nas duas versões com o mesmo protocolo.

\begin{figure}[H]
\centering
\begin{tikzpicture}[
  font=\small,
  % Same size for every box and fixed row positions, so both columns line up; flush center avoids
  % stretched spaces and hyphenation inside the boxes.
  box/.style={draw, rounded corners=2pt, text width=6.9cm, minimum height=1.15cm, align=flush center,
    inner sep=3pt, fill=#1},
  box/.default={white},
  arr/.style={-{Stealth[length=2.2mm]}, thick}]
  \def\colR{7.9}  % x of the right column
  \def\rowh{1.55} % distance between rows
  \node[box=gray!10] (dados) at (0,0) {\textbf{Spambase}: $\num{4601}\times57$ atributos,\\ rótulo spam/não spam};
  \node[box] (z) at (0,-\rowh) {Padronização z-score dos 57 atributos\\ (rótulos não usados)};
  \node[box=blue!8] (dkm) at (0,-2*\rowh) {\textbf{Questão 1}: double k-means, 9 pares $(K,H)$\\
    $\times$ 100 inicializações (sementes 0 a 99)};
  \node[box=blue!8] (best) at (0,-3*\rowh) {Melhor execução de cada par:\\ menor $W(\mathbf{G},\mathbf{U},\mathbf{V})$};
  \node[box=blue!8] (sil) at (0,-4*\rowh) {Silhueta euclidiana dos 9 vencedores;\\ $K^*=\arg\max_{(K,H)}\Sil$};
  \node[box=blue!8] (ystar) at (0,-5*\rowh) {Partição com $K^*$ grupos $\Rightarrow$ rótulos $y^*$;\\
    ARI frente ao rótulo original};

  \node[box=orange!10] (versoes) at (\colR,0) {\textbf{Questão 2}: versões $(\mathbf{X}, y)$ e $(\mathbf{X}, y^*)$\\
    da base, com os mesmos atributos brutos};
  \node[box=orange!10] (cv) at (\colR,-\rowh) {CV estratificada $30\times10$ (externa);\\
    ajuste interno por CV de 5 folds};
  \node[box=orange!10] (metr) at (\colR,-2*\rowh) {5 classificadores, 300 folds de teste:\\
    erro, precisão, cobertura e F-measure};
  \node[box=orange!10] (ic) at (\colR,-3*\rowh) {Estimativa pontual e IC 95\,\%\\ (média das 30 repetições)};
  \node[box=orange!10] (fried) at (\colR,-4*\rowh) {Teste de Friedman e pós-teste de Nemenyi\\ (300 blocos pareados)};
  \node[box=orange!10] (lc) at (\colR,-5*\rowh) {Curvas de aprendizagem (F-measure):\\
    treino de 5\,\% a 95\,\%, 10 divisões cada};

  \foreach \a/\b in {dados/z, z/dkm, dkm/best, best/sil, sil/ystar, versoes/cv, cv/metr, metr/ic, ic/fried, fried/lc}
    \draw[arr] (\a) -- (\b);
  \draw[arr] ($(dados.east)+(0,0.2)$) -- ($(versoes.west)+(0,0.2)$);
  \draw[arr] (ystar.east) -- ++(0.5,0) |- ($(versoes.west)+(0,-0.2)$);
\end{tikzpicture}
\caption{Encadeamento dos experimentos. A coluna da esquerda produz os rótulos $y^*$ usados na segunda versão da
base; a coluna da direita é executada igualmente para as duas versões. As curvas de aprendizagem usam divisões
próprias, independentes da validação cruzada $30\times10$.}\label{fig:pipeline}
\end{figure}

Os experimentos seguem cinco princípios, adotados para que a avaliação de cada modelo seja correta e a comparação
entre eles seja justa:

\begin{enumerate}
\item \textbf{Separação entre seleção e avaliação.} Os hiperparâmetros são escolhidos apenas com os dados de
  aprendizagem de cada fold externo, por uma validação cruzada interna; o fold de teste só é usado uma vez, para
  medir o modelo já ajustado. Escolher hiperparâmetros e medir o desempenho nos mesmos dados produz estimativas
  otimistas~\cite{varma2006bias,cawley2010overfitting}.
\item \textbf{Pré-processamento dentro do treino.} A padronização usada pelo $k$-NN e pela regressão logística faz
  parte do próprio modelo: média e desvio são estimados no conjunto em que o modelo é ajustado (cada fold interno
  na seleção, os 9 folds no reajuste final) e apenas aplicados ao conjunto de validação ou teste.
\item \textbf{Comparação pareada.} Os cinco classificadores são avaliados exatamente nas mesmas 300 partições
  externas e, dentro de cada uma, nas mesmas 5 partições internas. As diferenças entre eles refletem os modelos, e
  não sorteios diferentes, e cada fold externo forma um bloco para o teste de Friedman~\cite{friedman1937ranks} (ver a ressalva na Seção~\ref{sec:friedman}).
\item \textbf{Estratificação.} Todas as divisões (externas, internas e das curvas de aprendizagem) preservam as
  proporções das classes, o que reduz a variância das estimativas~\cite{kohavi1995study} e é indispensável na
  versão com $K^*$ grupos, cuja classe rara tem 34 exemplos.
\item \textbf{Reprodutibilidade.} Todas as fontes de aleatoriedade têm semente fixa (Tabela~\ref{tab:sementes}) e
  os resultados brutos de cada execução e de cada fold são salvos em \texttt{results/}.
\end{enumerate}

\begin{table}[H]
\centering
\caption{Fontes de aleatoriedade e sementes. No double k-means, cada uma das 100 execuções de um par usa uma
semente; na CV interna das curvas, $r=0,\dots,9$ é o índice da divisão.}\label{tab:sementes}
\small
\begin{tabularx}{\linewidth}{@{}l L c@{}}
\toprule
Etapa & Procedimento & Semente \\
\midrule
Q1, double k-means & Partições iniciais aleatórias de objetos e variáveis & 0 a 99 \\
Q2, CV externa & \texttt{RepeatedStratifiedKFold}, 10 folds $\times$ 30 repetições & 42 \\
Q2, CV interna & \texttt{StratifiedKFold}, 5 folds com embaralhamento & 42 \\
Q2, regressão logística & Semente do otimizador (\texttt{lbfgs} é determinístico) & 42 \\
Curvas de aprendizagem & \texttt{StratifiedShuffleSplit}, 10 divisões por proporção & 42 \\
Curvas, CV interna & \texttt{StratifiedKFold} com $\min(5, n_{\min})$ folds & $42+r$ \\
\bottomrule
\end{tabularx}
\end{table}

%======================================================================
\section{Questão 1: double k-means}\label{sec:q1}
%======================================================================

\subsection{O algoritmo}

O double k-means~\cite{vichi2001double} é um método de agrupamento simultâneo de objetos e variáveis
(\textit{co-clustering})~\cite{govaert2013coclustering}. Dada a matriz $\mathbf{X}=(x_{ij})$, $n\times d$, ele
obtém uma partição exclusiva dos objetos em $K$ grupos $P_1,\dots,P_K$ (matriz indicadora $\mathbf{U}=(u_{ik})$),
uma partição exclusiva das variáveis em $H$ grupos $Q_1,\dots,Q_H$ (matriz indicadora $\mathbf{V}=(v_{jh})$) e uma
matriz de protótipos $\mathbf{G}=(g_{kh})$, $K\times H$, um valor por bloco, que minimizam
\begin{equation}
W(\mathbf{G},\mathbf{U},\mathbf{V}) = \sum_{k=1}^{K}\sum_{h=1}^{H}\sum_{i=1}^{n}\sum_{j=1}^{d}
  u_{ik}\,v_{jh}\,(x_{ij}-g_{kh})^2 .
\label{eq:W}
\end{equation}
Partindo de partições aleatórias, o algoritmo alterna três passos até que nenhum objeto e nenhuma variável mude de grupo, como foi visto na disciplina:
\begin{enumerate}
\item com $\mathbf{U}$ e $\mathbf{V}$ fixos, $g_{kh}=\frac{1}{n_k n_h}\sum_{i}\sum_{j}u_{ik}v_{jh}x_{ij}$ (média do
  bloco $kh$);
\item com $\mathbf{G}$ e $\mathbf{V}$ fixos, cada objeto vai para o grupo $k$ que minimiza
  $\sum_h\sum_j v_{jh}(x_{ij}-g_{kh})^2$;
\item com $\mathbf{G}$ e $\mathbf{U}$ fixos, cada variável vai para o grupo $h$ que minimiza
  $\sum_k\sum_i u_{ik}(x_{ij}-g_{kh})^2$.
\end{enumerate}
Para o tratamento de empates e de grupos que fiquem vazios, adotou-se na implementação duas regras: (a) Um item só muda de grupo se o custo cair mais 
de $10^{-9}$ (melhora estrita, como em Hartigan e Wong~\cite{hartigan1979kmeans}), o que impede trocas causadas por arredondamento, e (b) um grupo esvaziado 
nos passos 2 ou 3 recebe o item de maior custo em relação ao próprio protótipo, escolhido entre os grupos com pelo menos dois membros, a mesma estratégia
do k-means do scikit-learn. Com essas regras $W$ nunca aumenta, toda execução termina com exatamente $K$ e $H$ grupos e nenhuma precisa ser descartada.
O limite de 100 iterações é só uma salvaguarda: as 900 execuções convergiram pelo critério dos slides.

\subsection{Protocolo}\label{sec:q1-protocolo}

\begin{itemize}
\item \textbf{Pré-processamento:} foi realizada a padronização z-score de cada atributo na base completa, pois sem a mesma, $W$ seria
  dominado pelos atributos \texttt{capital\_run\_length\_*}, cuja variância é milhares de vezes maior que a das frequências, e as variáveis 
  seriam agrupadas pela escala. O rótulo spam/não spam não entra no agrupamento, de forma que padronizar com a base completa não introduz vazamento.
\item \textbf{Grade:} $K\in\{2,3,4\}$ e $H\in\{1,\dots,K\}$, isto é, 9 pares: $(2,1)$, $(2,2)$, $(3,1)$, $(3,2)$,
  $(3,3)$, $(4,1)$, $(4,2)$, $(4,3)$ e $(4,4)$.
\item \textbf{Repetições:} 100 execuções por par a partir de inicializações aleatórias (sementes 0 a 99), 900 no
  total. Como o algoritmo só garante um mínimo local, a solução de cada par é a execução de menor $W$ (em caso
  de empate exato, a de menor semente).
\item \textbf{Escolha de $K^*$:} $W$ não serve para comparar pares diferentes, porque diminui quando $K$ e $H$ aumentam (há mais protótipos para ajustar). 
  Consequentemente, calculou-se a silhueta da partição de objetos de cada um dos 9 vencedores e escolheu-se $K^*$ pelo par de maior silhueta. 
  A silhueta de um objeto $i$ é $s(i)=\frac{b(i)-a(i)}{\max\{a(i),b(i)\}}$, em que $a(i)$ é a distância média de $i$ aos objetos do seu grupo e
  $b(i)$ a menor distância média de $i$ aos objetos de outro grupo; $\Sil$ é a média de $s(i)$ sobre os $n$ objetos~\cite{rousseeuw1987silhouettes}. 
  Usou-se a distância euclidiana nos 57 atributos padronizados, o mesmo espaço em que o algoritmo foi executado.
\item \textbf{Avaliação externa:} índice de Rand corrigido~\cite{hubert1985comparing} entre a partição de
  objetos com $K^*$ grupos e as classes a priori. Com $n_{ij}$ a tabela de contingência, $a_i$ e $b_j$ seus totais
  marginais,
  \begin{equation}
  \ARI=\frac{\sum_{ij}\binom{n_{ij}}{2}-\Big[\sum_i\binom{a_i}{2}\sum_j\binom{b_j}{2}\Big]\Big/\binom{n}{2}}
  {\frac12\Big[\sum_i\binom{a_i}{2}+\sum_j\binom{b_j}{2}\Big]-\Big[\sum_i\binom{a_i}{2}\sum_j\binom{b_j}{2}\Big]\Big/\binom{n}{2}} .
  \end{equation}
  O ARI vale 1 para partições idênticas e tem valor esperado 0 quando as partições são independentes; valores
  negativos indicam concordância abaixo da esperada ao acaso.
\end{itemize}

\subsection{Seleção de \texorpdfstring{$(K,H)$}{(K,H)}}

A Tabela~\ref{tab:q1grade} traz, para cada par, a melhor execução e sua silhueta, e a
Figura~\ref{fig:silhueta} mostra o gráfico $\Sil\times(K,H)$ pedido no enunciado.

\begin{table}[H]
\centering
\caption{Melhor das 100 execuções de cada par $(K,H)$ e silhueta da sua partição de objetos; todas as 900
execuções convergiram.}\label{tab:q1grade}
\small
\begin{tabular}{@{}cc S[table-format=6.2] c c S[table-format=1.4]@{}}
\toprule
$K$ & $H$ & {Menor $W$} & Semente & Iterações & {$\Sil$} \\
\midrule
2 & 1 & 256513.75 & 0 & 11 & 0.2961 \\
\textbf{2} & \textbf{2} & \bfseries 244486.70 & \textbf{0} & \textbf{31} & \bfseries 0.6596 \\
3 & 1 & 254414.01 & 0 & 40 & 0.2515 \\
3 & 2 & 239421.79 & 82 & 40 & 0.2367 \\
3 & 3 & 234585.87 & 7 & 30 & 0.1409 \\
4 & 1 & 253278.68 & 0 & 38 & 0.0484 \\
4 & 2 & 233358.70 & 3 & 43 & 0.2428 \\
4 & 3 & 228516.44 & 8 & 26 & 0.1620 \\
4 & 4 & 227250.88 & 18 & 28 & 0.1775 \\
\bottomrule
\end{tabular}
\end{table}

\begin{figure}[H]
\centering
\begin{tikzpicture}
\begin{axis}[
  width=0.92\linewidth, height=5.6cm,
  ybar, bar width=16pt, bar shift=0pt,
  ymin=0, ymax=0.75,
  xtick={1,...,9},
  xticklabels={{(2,1)},{(2,2)},{(3,1)},{(3,2)},{(3,3)},{(4,1)},{(4,2)},{(4,3)},{(4,4)}},
  xlabel={Par $(K,H)$}, ylabel={Silhueta},
  yticklabel style={/pgf/number format/.cd, fixed, fixed zerofill, precision=1, use comma},
  ymajorgrids, grid style={dashed, gray!30},
  nodes near coords,
  nodes near coords style={font=\scriptsize, /pgf/number format/.cd, fixed, fixed zerofill, precision=4, use comma},
  enlarge x limits=0.07,
]
\addplot[fill=blue!35, draw=blue!60!black] coordinates
  {(1,0.2961) (3,0.2515) (4,0.2367) (5,0.1409) (6,0.0484) (7,0.2428) (8,0.1620) (9,0.1775)};
\addplot[fill=orange!70, draw=orange!60!black] coordinates {(2,0.6596)};
\end{axis}
\end{tikzpicture}
\caption{Silhueta da melhor partição (menor $W$) de cada par $(K,H)$. O máximo ocorre em $(2,2)$, logo $K^*=2$
(e $H^*=2$).}\label{fig:silhueta}
\end{figure}

O par $(2,2)$ tem silhueta $\num{0.6596}$, mais que o dobro da segunda maior ($\num{0.2961}$, em $(2,1)$), e foi
selecionado: $K^*=2$. Na escala de interpretação de Kaufman e Rousseeuw~\cite{kaufman1990finding}, $\num{0.66}$ indica uma estrutura
razoável, enquanto os demais pares, entre $\num{0.05}$ e $\num{0.30}$, ficam nas faixas de estrutura fraca ou
inexistente. Note também que, como esperado, o menor $W$ cai quando $K$ e $H$ crescem, de $\num{256513.75}$ em $(2,1)$ a $\num{227250.88}$ em $(4,4)$.

As partições dos outros pares ajudam a entender a escolha: com $K\ge3$, o grupo dos 34 e-mails descrito na Seção~\ref{sec:particao} reaparece, inteiro 
em $(3,2)$ e $(3,3)$ e quase inteiro nos demais, e o restante da base é dividido em dois ou três grupos grandes (por exemplo, $\num{3437}$ e $\num{1130}$ objetos em $(3,2)$). 
Essa divisão da massa principal não tem separação nítida e derruba a silhueta, que fica entre $\num{0.05}$ e $\num{0.25}$ para $K\ge3$. 
Já em $(2,1)$, com um único grupo de variáveis, os 34 e-mails ficam diluídos num grupo de $\num{1018}$ objetos, e a silhueta não passa de $\num{0.30}$.

\subsection{Partição selecionada}\label{sec:particao}

A melhor execução de $(K^*,H^*)=(2,2)$ , com semente 0, separa as variáveis em um grupo $Q_1$ com 11 atributos, \textit{hp}, \textit{hpl}, \textit{650}, 
\textit{lab}, \textit{labs}, \textit{telnet}, \textit{857}, \textit{415}, \textit{85}, \textit{technology} e \textit{direct}, e um grupo $Q_2$ com os 
outros 46 (37 palavras, os 6 caracteres e os 3 atributos de maiúsculas). Os 11 atributos de $Q_1$ são nomes e números ligados aos laboratórios da HP 
em Palo Alto (códigos de área 650 e 415), a origem dos e-mails legítimos da base. Os objetos ficam em um grupo $P_1$ com $\num{4567}$ e-mails e um 
grupo $P_2$ com 34,\footnote{No código e nos arquivos de resultados, os grupos $P_1,P_2$ e $Q_1,Q_2$ aparecem com os rótulos 0 e 1.} cujos protótipos, 
em unidades de desvio-padrão, formam a matriz $\mathbf{G}$ da Tabela~\ref{tab:G}.

\begin{table}[H]
\centering
\caption{Matriz de protótipos dos blocos $\mathbf{G}$ para $(K^*,H^*)=(2,2)$ (atributos padronizados).}\label{tab:G}
\small
\begin{tabular}{@{}l S[table-format=-1.4] S[table-format=-1.4] r@{}}
\toprule
 & {$Q_1$ (11 variáveis HP)} & {$Q_2$ (46 variáveis)} & Objetos \\
\midrule
$P_1$ & -0.0511 & 0.0011 & \num{4567} \\
$P_2$ & 6.8612 & -0.1444 & 34 \\
\bottomrule
\end{tabular}
\end{table}

O bloco $P_2\times Q_1$ tem protótipo $\num{6.86}$: os 34 e-mails de $P_2$ usam as 11 palavras da HP, em média, quase sete desvios-padrão acima da 
média da base, enquanto nas demais variáveis ($Q_2$) ficam ligeiramente abaixo dela ($\num{-0.14}$). O grupo $P_1$, por sua vez, tem protótipos
próximos de zero nos dois blocos, ou seja, reúne o restante da base sem distinção.

\begin{table}[H]
\centering
\caption{Matriz de confusão (contingência) entre a partição de objetos do double k-means e as classes a priori.}\label{tab:confusao}
\small
\begin{tabular}{@{}l r r r@{}}
\toprule
Classe a priori & $P_1$ & $P_2$ & Total \\
\midrule
Não spam (0) & \num{2754} & 34 & \num{2788} \\
Spam (1) & \num{1813} & 0 & \num{1813} \\
\midrule
Total & \num{4567} & 34 & \num{4601} \\
\bottomrule
\end{tabular}
\end{table}

Quando essa partição é cruzada com as classes a priori (Tabela~\ref{tab:confusao}), os 34 e-mails de $P_2$ são todos não spam e $P_1$ 
reúne os $\num{1813}$ spam e os outros $\num{2754}$ não spam, o que resulta em um índice de Rand corrigido de $\ARI=\num{-0.00495}$ 
entre as duas partições.

\todoequipe{Comentário sobre o ARI com Bruna e Everthon. Interpretar o valor $\ARI=\num{-0.00495}$ à luz da Tabela~\ref{tab:confusao} e da 
matriz $\mathbf{G}$: a.~o que significa um ARI próximo de zero e negativo; b.~por que ele é compatível com uma silhueta alta; 
c.~o que a partição encontrada revela sobre a estrutura da base.}

Um ARI tão próximo de zero indica que a partição obtida não guarda relação com as classes spam/não spam, pois a concordância entre as duas é praticamente
a que se esperaria de partições independentes: o único corte feito pelo double k-means passa por dentro da classe não spam, deixando todo o spam 
misturado ao restante em $P_1$, de modo que a pureza da partição (60,6\,\%) é exatamente a proporção da classe majoritária. O sinal negativo decorre da
composição de $P_2$, porque um grupo de 34 e-mails sorteados ao acaso teria, em média, cerca de 21 não spam e 13 spam, e separá-lo do restante 
desfaria cerca de 81~mil pares de e-mails da mesma classe, ao passo que $P_2$, formado apenas por não spam, desfaz $34\times\num{2754}=\num{93636}$ 
desses pares, o que deixa a concordância com as classes ligeiramente abaixo da esperada ao acaso. 

Esse valor é compatível com a silhueta alta porque os dois índices medem propriedades diferentes: a silhueta é um critério interno, calculado apenas 
com as distâncias entre os objetos, e é elevada porque $P_2$ está, em média, $\num{6.86}$ desvios-padrão acima da média nas 11 variáveis de $Q_1$
(Tabela~\ref{tab:G}), o que torna $s(i)$ alto para quase todos os objetos de $P_1$ (Seção~\ref{sec:q1-discussao}), enquanto o ARI é um critério externo, 
que compara a partição com um rótulo que o algoritmo nunca utilizou. A partição revela, portanto, que a estrutura geométrica dominante na base 
padronizada não é a divisão entre spam e não spam, mas um pequeno grupo de e-mails internos da HP com uso intenso de termos como \textit{hp}, 
\textit{hpl}, \textit{650} e \textit{lab}, reflexo da forma como os e-mails legítimos da base foram coletados (Seção~\ref{sec:dados}). A diferença entre 
as classes está em direções mais sutis do espaço de atributos, que o agrupamento não captura, mas que os classificadores supervisionados da Questão~2 
exploram, com F-measure macro de $\num{0.923}$ na regressão logística (Tabela~\ref{tab:ic-original}).

A Figura~\ref{fig:convergencia} mostra a função objetivo ao longo das iterações da execução selecionada. $W$ parte
de $\num{262251.38}$ na partição inicial e decresce monotonamente até $\num{244486.70}$, quando nenhum objeto ou
variável muda de grupo (iteração 31, que repete o valor da 30). Entre as iterações 11 e 20 há um patamar em que
$W$ quase não varia, seguido de uma nova queda até a convergência: o algoritmo pode passar várias iterações com
transferências de pouco ganho antes de encontrar uma reorganização que reduz $W$ de forma apreciável.

\begin{figure}[H]
\centering
\begin{tikzpicture}
\begin{axis}[
  width=0.9\linewidth, height=5.6cm,
  xlabel={Iteração (0 = partição inicial)},
  ylabel={$W(\mathbf{G},\mathbf{U},\mathbf{V})$},
  scaled y ticks=false,
  yticklabel style={/pgf/number format/.cd, fixed, use comma, 1000 sep={.}},
  xmin=-1, xmax=32,
  grid=major, grid style={dashed, gray!30},
]
\addplot[blue!70!black, mark=*, mark size=1.3pt]
  table[x=iteration, y=objective, col sep=comma] {../../results/q1_objective_history.csv};
\end{axis}
\end{tikzpicture}
\caption{Função objetivo versus iterações na melhor execução de $(K^*,H^*)=(2,2)$. Cada ponto é $W$ com a
partição do fim da iteração e os protótipos ótimos para ela.}\label{fig:convergencia}
\end{figure}

\subsection{Discussão}\label{sec:q1-discussao}

A silhueta escolheu uma partição muito desbalanceada porque, sendo $\Sil$ a média de $s(i)$ sobre todos os
objetos, 99,26\,\% dos termos dessa média vêm de objetos de $P_1$. Para cada um deles, $a(i)$ é a distância média ao grosso da base,
e $b(i)$ é a distância média aos 34 objetos de $P_2$, que estão, em média, quase sete desvios-padrão acima da média
em 11 variáveis (Tabela~\ref{tab:G}); assim, $b(i)$ tende a ser bem maior que $a(i)$ e $s(i)$ fica alto. A
silhueta premia, portanto, a separação de um pequeno grupo atípico, e não uma divisão equilibrada da base. Esse
comportamento é coerente com duas limitações do double k-means apontadas nos slides da disciplina, a sensibilidade
a valores discrepantes e a preferência por grupos compactos e esféricos, e com o fato de a soma de quadrados dar
peso desproporcional às caudas longas dos atributos esparsos (Seção~\ref{sec:dados}). Para a Questão~2, a
consequência é que a segunda versão da base pede a distinção entre esses 34 e-mails e o restante, uma tarefa
diferente e muito mais desbalanceada que a detecção de spam.

Essa tarefa tem, além disso, uma fronteira linear conhecida, porque, na convergência, cada objeto está no grupo de
menor custo do passo~2, dados $\mathbf{G}$ e $\mathbf{V}$. Com $\mathbf{z}_i$ o vetor padronizado do objeto $i$, $c_k(\mathbf{z}_i)$ o custo
de colocá-lo no grupo $k$ e $d_h$ o número de variáveis de $Q_h$, os termos $z_{ij}^2$ se cancelam na diferença
entre os custos dos dois grupos, que é linear em $\mathbf{z}_i$:
\begin{equation}
c_1(\mathbf{z}_i)-c_2(\mathbf{z}_i)=\sum_{h=1}^{H}\Big[2\,(g_{2h}-g_{1h})\,S_{ih}-d_h\,\big(g_{2h}^2-g_{1h}^2\big)\Big],
\qquad S_{ih}=\sum_{j\in Q_h}z_{ij}.
\label{eq:fronteira}
\end{equation}
Com os valores da Tabela~\ref{tab:G}, o objeto vai para $P_2$ se e somente se $S_{i1}-\num{0.021}\,S_{i2}>\num{37.5}$,
isto é, quando a soma das 11 variáveis da HP padronizadas passa de cerca de 37,5 (em média, mais de 3,4
desvios-padrão por variável). Aplicada aos dados padronizados, essa regra reproduz exatamente os $\num{4601}$
rótulos, de modo que as duas classes da versão \texttt{kstar\_classes} são linearmente separáveis no espaço
padronizado e também no espaço original, já que a padronização é uma transformação afim.

%======================================================================
\section{Questão 2: classificação supervisionada}\label{sec:q2}
%======================================================================

\subsection{Versões da base}

\begin{table}[H]
\centering
\caption{Versões da base usadas na Questão 2. Os atributos são os 57 originais, sem transformação prévia.}\label{tab:versoes}
\small
\begin{tabular}{@{}l l r r@{}}
\toprule
Versão & Variável resposta & Classe 0 & Classe 1 \\
\midrule
\texttt{original\_2\_classes} & spam/não spam (a priori) & \num{2788} (60,6\,\%) & \num{1813} (39,4\,\%) \\
\texttt{kstar\_classes} & grupos $P_1$/$P_2$ da Questão~1 ($K^*=2$) & \num{4567} (99,26\,\%) & 34 (0,74\,\%) \\
\bottomrule
\end{tabular}
\end{table}

Como $K^*=2$, a segunda versão também é binária, e a sua classe rara tem apenas 3 ou 4 exemplos nos folds externos
de teste ($\approx$\,460 exemplos) e cerca de 30 nos folds de aprendizagem ($\approx$\,\num{4141} exemplos).

\subsection{Classificadores}\label{sec:classificadores}

Todos os classificadores seguem a interface do scikit-learn~\cite{pedregosa2011scikit}; os quatro primeiros correspondem às regras pedidas no
enunciado, e o quinto combina os quatro.

\paragraph{i) Bayesiano gaussiano.} Atribui $\mathbf{x}$ à classe $\omega_l$ de maior posteriori
$P(\omega_l\mid\mathbf{x})\propto p(\mathbf{x}\mid\omega_l)\,P(\omega_l)$, com estimativas de máxima verossimilhança~\cite{duda2001pattern}
\begin{equation}
\hat P(\omega_i)=\frac{n_i}{n},\qquad
\hat{\bm\mu}_i=\frac{1}{n_i}\sum_{\mathbf{x}_k\in\omega_i}\mathbf{x}_k,\qquad
\hat{\bm\Sigma}_i=\frac{1}{n_i}\sum_{\mathbf{x}_k\in\omega_i}(\mathbf{x}_k-\hat{\bm\mu}_i)(\mathbf{x}_k-\hat{\bm\mu}_i)^\top ,
\end{equation}
e densidade normal multivariada. A decisão é tomada em escala logarítmica,
\begin{equation}
\hat\omega(\mathbf{x})=\arg\max_i\Big[\ln\hat P(\omega_i)-\tfrac12\ln\big|\tilde{\bm\Sigma}_i\big|
-\tfrac12(\mathbf{x}-\hat{\bm\mu}_i)^\top\tilde{\bm\Sigma}_i^{-1}(\mathbf{x}-\hat{\bm\mu}_i)\Big],
\qquad \tilde{\bm\Sigma}_i=\hat{\bm\Sigma}_i+\varepsilon\mathbf{I},
\end{equation}
com o determinante e a forma quadrática calculados pela fatoração de Cholesky de $\tilde{\bm\Sigma}_i$. O termo
$\varepsilon\mathbf{I}$, com $\varepsilon=10^{-6}$ fixo, foi necessário porque a estimativa de máxima
verossimilhança é singular nas duas versões da base. Em \texttt{kstar\_classes}, a classe rara tem cerca de 30
exemplos de treino para $d=57$ atributos, e $\operatorname{posto}(\hat{\bm\Sigma}_i)\le n_i-1<d$; além disso, os 34
e-mails de $P_2$ são todos não spam, só 28 deles são distintos e 29 dos 57 atributos são nulos em todos eles, de
modo que a covariância estimada com os 34 tem posto 16. Na base original, \texttt{word\_freq\_cs} é não nulo em
apenas 1 dos $\num{1813}$ e-mails spam, e \texttt{word\_freq\_857} em apenas 2. Sempre que esses exemplos ficam fora
do treino, o atributo é constante na classe spam e $\hat{\bm\Sigma}_i$ é singular. Isso ocorre em pelo menos 30 dos
300 folds externos (em cada repetição, o fold cujo teste contém o único spam com \texttt{word\_freq\_cs} não nulo),
em parte dos folds internos e na maior parte das divisões com treino pequeno da curva de aprendizagem, casos em
que, sem a correção, a densidade não estaria definida. O valor $\varepsilon=10^{-6}$ é o padrão do parâmetro
\texttt{reg\_covar} do \texttt{GaussianMixture} do scikit-learn e é muito menor que a menor variância dos
atributos ($\num{0.076}^2\approx\num{0.0058}$), de modo que só altera de forma apreciável as direções de variância
nula. Trata-se de uma forma mínima da regularização por encolhimento da
covariância~\cite{ledoit2004wellconditioned,hastie2009elements} e, como $\varepsilon$ não é ajustado, o
classificador continua sem hiperparâmetros.

\paragraph{ii) Bayesiano baseado em $k$ vizinhos.} Estimando $p(\mathbf{x}\mid\omega_i)=k_i/(n_iV)$, em que $V$ é
o volume da menor vizinhança de $\mathbf{x}$ que contém $k$ exemplos de treino e $k_i$ deles são da classe
$\omega_i$, e $\hat P(\omega_i)=n_i/n$, a regra de Bayes se reduz a $P(\omega_i\mid\mathbf{x})=k_i/k$, isto é,
ao voto majoritário entre os $k$ vizinhos~\cite{duda2001pattern,cover1967nearest}. A vizinhança é definida
pela distância euclidiana, city-block ($L_1$) ou Chebyshev ($L_\infty$), calculada sobre os atributos
padronizados (z-score estimado no conjunto de treino), e tanto essa distância quanto o número de vizinhos $k$ são
hiperparâmetros.

\paragraph{iii) Bayesiano com janela de Parzen.} A densidade de cada classe é estimada pelo kernel multivariado
produto com kernel gaussiano univariado e a mesma janela $h$ em todas as dimensões~\cite{parzen1962estimation,duda2001pattern},
\begin{equation}
\hat p(\mathbf{x}\mid\omega_i)=\frac{1}{n_i}\sum_{\mathbf{x}_m\in\omega_i}\prod_{j=1}^{d}\frac{1}{h}\,
\phi\!\left(\frac{x_j-x_{mj}}{h}\right)
=\frac{1}{n_i(2\pi)^{d/2}h^d}\sum_{\mathbf{x}_m\in\omega_i}\exp\!\left(-\frac{\lVert\mathbf{x}-\mathbf{x}_m\rVert^2}{2h^2}\right),
\end{equation}
com $\phi$ a densidade normal padrão e $\hat P(\omega_i)=n_i/n$; a decisão maximiza
$\ln\hat p(\mathbf{x}\mid\omega_i)+\ln\hat P(\omega_i)$. A janela $h$ é o único hiperparâmetro e, nesta
implementação, o estimador recebe os atributos na escala original (Seção~\ref{sec:discussao}).

\paragraph{iv) Regressão logística.} Modelo discriminativo
$P(\omega_1\mid\mathbf{x})=\sigma(\mathbf{w}^\top\mathbf{z}+b)$, em que $\mathbf{z}$ são os atributos padronizados
no treino e $\sigma$ é a função logística, ajustado por máxima verossimilhança penalizada~\cite{hastie2009elements}:
$\min_{\mathbf{w},b}\ \tfrac12\lVert\mathbf{w}\rVert^2+C\sum_k\ell(y_k,\mathbf{x}_k)$, com $\ell$ a perda logística,
otimizador L-BFGS e até $\num{10000}$ iterações~\cite{pedregosa2011scikit}, tendo como hiperparâmetro o
inverso da força de regularização, $C$.

\paragraph{v) Voto majoritário.} Combina os rótulos preditos pelos quatro classificadores anteriores,
$\hat\omega(\mathbf{x})=\arg\max_c\sum_{m=1}^{4}\mathbf{1}[\hat\omega_m(\mathbf{x})=c]$~\cite{kuncheva2014combining}.
Cada membro entra com os hiperparâmetros que ele mesmo selecionou na validação cruzada interna do fold externo
corrente; o comitê, portanto, não tem hiperparâmetros próprios nem uma busca adicional. Com quatro votantes e duas
classes pode haver empate de 2 a 2, resolvido em favor do menor rótulo (classe 0), o comportamento do
\texttt{VotingClassifier} do scikit-learn.

\subsection{Protocolo de avaliação: validação cruzada \texorpdfstring{$30\times10$}{30×10} com ajuste interno}\label{sec:protocolo}

A Figura~\ref{fig:cvaninhada} ilustra o protocolo, executado da mesma forma para as duas versões da base:

\begin{enumerate}
\item A base é dividida em 10 folds estratificados; o procedimento é repetido 30 vezes com permutações diferentes
  (\texttt{RepeatedStratifiedKFold}, semente 42), totalizando 300 pares (aprendizagem, teste), com $\num{4140}$ ou
  $\num{4141}$ exemplos de aprendizagem e 460 ou 461 de teste.
\item Em cada par, o fold de teste é posto de lado. Os 9 folds de aprendizagem são divididos em 5 folds
  estratificados (\texttt{StratifiedKFold}, semente 42), e cada combinação de hiperparâmetros da grade é
  avaliada pela F-measure macro média nesses 5 folds (\texttt{GridSearchCV}). Escolhe-se a combinação de maior
  média; em caso de empate, a primeira da grade.
\item O modelo é treinado novamente nos 9 folds de aprendizagem com a combinação escolhida e avaliado uma única
  vez no fold de teste, nas quatro métricas.
\item O voto majoritário é montado com os quatro modelos do passo 3, retreinado nos 9 folds e avaliado no mesmo
  fold de teste.
\end{enumerate}

\begin{figure}[H]
\centering
\begin{tikzpicture}[font=\small]
  \foreach \i in {0,...,8} {
    \draw[fill=blue!12] (\i*1.45,0) rectangle ++(1.35,0.6);
    \node at (\i*1.45+0.675,0.3) {\footnotesize F\the\numexpr\i+1\relax};
  }
  \draw[fill=orange!55] (9*1.45,0) rectangle ++(1.35,0.6);
  \node at (9*1.45+0.675,0.3) {\footnotesize Teste};
  \node[anchor=south west] at (0,0.7) {\textbf{CV externa}: 10 folds estratificados, repetida 30 vezes (300 folds de teste)};
  \draw[decorate, decoration={brace, mirror, amplitude=5pt}] (0,-0.1) -- (8*1.45+1.35,-0.1)
    node[midway, below=6pt] {9 folds de aprendizagem ($\approx\num{4141}$ exemplos)};
  \foreach \j in {0,...,4} {
    \pgfmathsetmacro{\x}{\j*2.61}
    \ifnum\j=4
      \draw[fill=green!35] (\x,-2.2) rectangle ++(2.51,0.6);
      \node at (\x+1.255,-1.9) {\footnotesize validação};
    \else
      \draw[fill=green!10] (\x,-2.2) rectangle ++(2.51,0.6);
      \node at (\x+1.255,-1.9) {\footnotesize treino};
    \fi
  }
  \node[anchor=north west, text width=10.4cm] at (0,-2.35) {\textbf{CV interna}: 5 folds estratificados dos 9 folds
    de aprendizagem; cada fold é usado uma vez como validação. Seleciona $\theta^*$ pela \mbox{F-measure} macro média.};
  \draw[-{Stealth[length=2.2mm]}, thick] (12.95,-1.9) -- (13.73,-1.9) -- (13.73,-0.1);
  \node[anchor=north, text width=2.6cm, align=center, font=\footnotesize] at (13.6,-2.35)
    {reajuste com $\theta^*$ nos 9 folds e avaliação no teste};
\end{tikzpicture}
\caption{Validação cruzada com ajuste interno de hiperparâmetros. O fold de teste não participa da escolha de $\theta^*$ nem do ajuste da padronização.}\label{fig:cvaninhada}
\end{figure}

Por fold externo são 45 combinações de hiperparâmetros (1 do gaussiano, 33 do $k$-NN, 6 da janela de Parzen e 5 da regressão logística) avaliadas 
em 5 folds internos, isto é, $\num{67500}$ ajustes internos por versão da base, além dos reajustes. Como as mesmas partições externas e internas 
valem para todos os classificadores, a comparação é pareada, e o ajuste interno evita o viés de seleção que surgiria se se reportasse a melhor média 
da busca como estimativa de desempenho~\cite{varma2006bias,cawley2010overfitting}; a Tabela~\ref{tab:hiper} mostra que, na versão \texttt{kstar\_classes},
essa média interna chega a superar a F-measure de teste em mais de um ponto percentual.

\subsection{Métricas}

Em cada fold de teste calculam-se a taxa de erro ($1-$acurácia) e as versões macro da precisão, da cobertura (\textit{recall}) e da F-measure ($F_1$), 
isto é, a média simples dos valores de cada classe~\cite{sokolova2009systematic}. 
A média macro foi escolhida porque dá o mesmo peso às duas classes: em \texttt{kstar\_classes}, a taxa de erro e as médias ponderadas seriam 
dominadas pelo grupo $P_1$ e quase não distinguiriam um bom classificador do que sempre prediz $P_1$ (erro de 0,74\,\%, F-measure macro de
aproximadamente 0,50), razão pela qual a F-measure macro também foi adotada como critério de seleção de hiperparâmetros.

\subsection{Obtenção dos hiperparâmetros}\label{sec:hiper}

A Tabela~\ref{tab:grades} lista as grades: a do $k$-NN inclui valores ímpares de $k$ (sem empate de votos com duas classes) espaçados de forma 
aproximadamente geométrica e as três distâncias pedidas no enunciado, enquanto as de $h$ e $C$ cobrem várias ordens de grandeza. 
O código descarta os valores de $k$ maiores que o menor conjunto de treino interno, mas isso não chegou a ocorrer: mesmo no menor caso 
(115 exemplos, nas curvas de aprendizagem com 5\,\% da versão \texttt{kstar\_classes}), todos os valores da grade são válidos.

\begin{table}[H]
\centering
\caption{Grades de hiperparâmetros. A seleção usa a F-measure macro média na CV interna de 5 folds.}\label{tab:grades}
\small
\begin{tabularx}{\linewidth}{@{}l L c@{}}
\toprule
Classificador & Hiperparâmetros e valores candidatos & Combinações \\
\midrule
Bayesiano gaussiano & nenhum (estimativas de máxima verossimilhança; $\varepsilon=10^{-6}$ fixo) & 1 \\
Bayesiano $k$-NN & $k\in\{1,3,5,7,11,15,21,31,45,63,91\}$; distância $\in\{$euclidiana, city-block, Chebyshev$\}$ & 33 \\
Janela de Parzen & $h\in\{0{,}1;\,0{,}5;\,1;\,2;\,5;\,10\}$ & 6 \\
Regressão logística & $C\in\{0{,}01;\,0{,}1;\,1;\,10;\,100\}$ (penalidade $L_2$, L-BFGS) & 5 \\
Voto majoritário & nenhum (usa os membros já ajustados no mesmo fold) & -- \\
\bottomrule
\end{tabularx}
\end{table}

Como a busca é refeita em cada um dos 300 folds externos, cada classificador tem uma distribuição de valores
escolhidos, resumida na Tabela~\ref{tab:hiper} junto com a F-measure média da melhor combinação na CV interna.

\begin{table}[H]
\centering
\caption{Hiperparâmetros selecionados nos 300 folds externos (frequência entre parênteses) e F-measure macro
média da combinação escolhida na CV interna. Para comparação, a F-measure média de teste está entre colchetes.}\label{tab:hiper}
\small
\begin{tabularx}{\linewidth}{@{}l L L@{}}
\toprule
Classificador & \texttt{original\_2\_classes} & \texttt{kstar\_classes} \\
\midrule
Bayesiano $k$-NN &
  city-block: $k=5$~(147), $k=7$~(68), $k=1$~(56), $k=3$~(10); euclidiana: $k=1$~(17), $k=3$~(2).
  Interna \num{0.9086} [teste \num{0.9093}] &
  euclidiana: $k=1$~(191), $k=11$~(16), $k=3$~(15), outros (8); city-block: $k=1$~(51), $k=11$~(12), outros (7).
  Interna \num{0.9879} [teste \num{0.9722}] \\
\addlinespace
Janela de Parzen &
  $h=1$~(205), $h=2$~(95). Interna \num{0.8108} [teste \num{0.8221}] &
  $h=2$~(148), $h=1$~(115), $h=0{,}5$~(19), $h=0{,}1$~(18). Interna \num{0.9504} [teste \num{0.9545}] \\
\addlinespace
Regressão logística &
  $C=100$~(160), $C=10$~(126), $C=1$~(14). Interna \num{0.9226} [teste \num{0.9230}] &
  $C=1$~(163), $C=0{,}1$~(126), $C=10$~(10), $C=100$~(1). Interna \num{0.9803} [teste \num{0.9762}] \\
\addlinespace
Bayesiano gaussiano &
  sem hiperparâmetros. Interna \num{0.8283} [teste \num{0.8284}] &
  sem hiperparâmetros. Interna \num{0.8971} [teste \num{0.8952}] \\
\bottomrule
\end{tabularx}
\end{table}

Algumas observações sobre as escolhas:
\begin{itemize}
\item \textbf{$k$-NN.} Na base original, a distância city-block foi escolhida em 281 dos 300 folds, com $k$ entre 5 e 7 na maioria. 
  A distância $L_1$ soma as diferenças em todas as coordenadas e é menos afetada que a euclidiana por uma única coordenada extrema, o que convém 
  a atributos esparsos de cauda longa. Na versão   \texttt{kstar\_classes}, predominam $k=1$ (242 folds) e a distância euclidiana (230): com apenas 
  cerca de 30 exemplos da classe rara no treino, vizinhanças maiores passam a ser dominadas pela classe majoritária, e a distância euclidiana sobre 
  atributos padronizados é a mesma geometria em que o double k-means definiu os rótulos $y^*$ (Seção~\ref{sec:q1-discussao}). A   distância de Chebyshev, 
  que considera só a maior diferença entre coordenadas, não foi escolhida em nenhum fold.
\item \textbf{Janela de Parzen.} As escolhas ficam no meio da grade ($h\in\{1,2\}$ em todos os folds da base original e em 263 da 
  versão \texttt{kstar\_classes}), longe dos extremos: janelas pequenas tornam a densidade concentrada nos pontos de treino e janelas
  grandes a tornam quase uniforme.
\item \textbf{Regressão logística.} Na base original, a busca prefere regularização fraca ($C=100$ em 160 folds, $C=10$ em 126), e, como $C=100$ é 
  o limite superior da grade, valores maiores poderiam ter sido testados. Na versão \texttt{kstar\_classes}, a regularização é mais 
  forte ($C\le1$ em 289 folds), o que se explica pela separabilidade linear das duas classes dessa versão (Equação~\ref{eq:fronteira}): 
  em dados separáveis, a estimativa sem penalização não existe, porque os coeficientes crescem sem limite. Nesse caso é a penalização que define 
  o hiperplano estimado, e a validação interna preferiu penalizações mais fortes que na base original.
\item \textbf{Viés da seleção.} Na base original, a F-measure interna da combinação escolhida é praticamente igual à de teste, enquanto, 
  na versão \texttt{kstar\_classes}, a média interna do $k$-NN (\num{0.9879}) supera a de teste   (\num{0.9722}), porque escolher o máximo 
  entre 33 candidatos avaliados com poucos exemplos da classe rara favorece combinações com sorte na validação. É exatamente o otimismo que o 
  protocolo da Seção~\ref{sec:protocolo} evita ao medir o desempenho num fold que não participou da seleção.
\end{itemize}

\subsection{Estimativas pontuais e intervalos de confiança}\label{sec:ic}

Os 300 valores de cada métrica não são observações independentes: os folds de uma mesma repetição compartilham exemplos de treino, e as 30 repetições
reutilizam os mesmos dados. Tratar os 300 valores como uma amostra independente subestimaria a variância e estreitaria artificialmente o 
intervalo~\cite{dietterich1998approximate,nadeau2003inference}. 
Por isso, para cada classificador e métrica, calculou-se primeiro a média $\bar m_r$ dos 10 folds de cada repetição $r$ 
(que é, ela própria, uma estimativa por validação cruzada de 10 folds) e, a partir das 30 médias, a estimativa pontual e o intervalo de 
confiança de 95\,\% pela distribuição $t$ de Student com 29 graus de liberdade:
\begin{equation}
\hat m=\frac{1}{30}\sum_{r=1}^{30}\bar m_r,\qquad
\hat m\pm t_{0{,}975;\,29}\,\frac{s_{\bar m}}{\sqrt{30}},\qquad t_{0{,}975;\,29}\approx\num{2.045},
\end{equation}
em que $s_{\bar m}$ é o desvio-padrão amostral das 30 médias. As Tabelas~\ref{tab:ic-original}
e~\ref{tab:ic-kstar} apresentam os resultados, e a Figura~\ref{fig:boxplots} a distribuição dos valores por fold.

\begin{table}[H]
\centering
\caption{\texttt{original\_2\_classes}: estimativa pontual $\pm$ meia-amplitude do IC 95\,\% (30 repetições), com
precisão, cobertura e F-measure em média macro e o melhor valor de cada métrica em negrito.}\label{tab:ic-original}
\small
\begin{tabular*}{\linewidth}{@{\extracolsep{\fill}}l c c c c@{}}
\toprule
Classificador & Taxa de erro & Precisão & Cobertura & F-measure \\
\midrule
Bayesiano gaussiano & \ci{0.1705}{0.0004} & \ci{0.8362}{0.0004} & \ci{0.8507}{0.0004} & \ci{0.8284}{0.0004} \\
Bayesiano $k$-NN & \ci{0.0855}{0.0009} & \ci{0.9163}{0.0009} & \ci{0.9046}{0.0010} & \ci{0.9093}{0.0009} \\
Janela de Parzen & \ci{0.1686}{0.0010} & \ci{0.8254}{0.0011} & \ci{0.8201}{0.0013} & \ci{0.8221}{0.0011} \\
Regressão logística & \ci{0.0730}{0.0004} & \ci{0.9265}{0.0005} & \best{\ci{0.9204}{0.0005}} & \ci{0.9230}{0.0005} \\
Voto majoritário & \best{\ci{0.0716}{0.0008}} & \best{\ci{0.9331}{0.0007}} & \ci{0.9174}{0.0009} & \best{\ci{0.9238}{0.0009}} \\
\bottomrule
\end{tabular*}
\end{table}

\begin{table}[H]
\centering
\caption{\texttt{kstar\_classes}: estimativa pontual $\pm$ meia-amplitude do IC 95\,\% (30 repetições), com
precisão, cobertura e F-measure em média macro, a taxa de erro com cinco casas decimais e o melhor valor de cada
métrica em negrito.}\label{tab:ic-kstar}
\small
\begin{tabular*}{\linewidth}{@{\extracolsep{\fill}}l c c c c@{}}
\toprule
Classificador & Taxa de erro & Precisão & Cobertura & F-measure \\
\midrule
Bayesiano gaussiano & \ci{0.00251}{0.00012} & \ci{0.9540}{0.0079} & \ci{0.8705}{0.0066} & \ci{0.8952}{0.0065} \\
Bayesiano $k$-NN & \ci{0.00078}{0.00005} & \ci{0.9799}{0.0018} & \ci{0.9729}{0.0030} & \ci{0.9722}{0.0023} \\
Janela de Parzen & \ci{0.00122}{0.00011} & \ci{0.9737}{0.0045} & \ci{0.9471}{0.0054} & \ci{0.9545}{0.0043} \\
Regressão logística & \best{\ci{0.00077}{0.00011}} & \ci{0.9641}{0.0044} & \best{\ci{0.9949}{0.0041}} & \best{\ci{0.9762}{0.0040}} \\
Voto majoritário & \ci{0.00080}{0.00009} & \best{\ci{0.9962}{0.0022}} & \ci{0.9490}{0.0045} & \ci{0.9668}{0.0034} \\
\bottomrule
\end{tabular*}
\end{table}

\begin{figure}[!tbp]
\centering
\begin{tikzpicture}
\begin{groupplot}[
  group style={group size=2 by 4, horizontal sep=1.7cm, vertical sep=0.7cm, xticklabels at=edge bottom},
  scale only axis, width=6.2cm, height=3.4cm,
  boxplot/draw direction=y, boxplot/box extend=0.55,
  % Mean as a white diamond: the default marker inherits the box fill and disappears inside the box.
  boxplot/draw/average/.code={\draw[black] \pgfextra \pgfsetfillcolor{white}%
    \pgftransformshift{\pgfplotsboxplotpointabbox{\pgfplotsboxplotvalue{average}}{0.5}}%
    \pgfuseplotmark{diamond*}\endpgfextra;},
  xmin=0.4, xmax=5.6, xtick={1,...,5}, xticklabels={GB,KNN,PZ,RL,VM},
  ymajorgrids, grid style={dashed, gray!30},
  scaled y ticks=false,
  yticklabel style={/pgf/number format/.cd, fixed, precision=3, use comma},
  tick label style={font=\footnotesize}, label style={font=\footnotesize}, title style={font=\small},
]
\nextgroupplot[title={\texttt{original\_2\_classes}}, ylabel={Taxa de erro}]\bppanel{0}
\nextgroupplot[title={\texttt{kstar\_classes}}]\bppanel{5}
\nextgroupplot[ylabel={Precisão (macro)}]\bppanel{10}
\nextgroupplot\bppanel{15}
\nextgroupplot[ylabel={Cobertura (macro)}]\bppanel{20}
\nextgroupplot\bppanel{25}
\nextgroupplot[ylabel={F-measure (macro)}]\bppanel{30}
\nextgroupplot\bppanel{35}
\end{groupplot}
\end{tikzpicture}
\caption{Distribuição das quatro métricas nos 300 folds de teste (colunas: versões da base; linhas: métricas;
cada painel tem escala própria). A caixa vai do primeiro ao terceiro quartil, com a mediana; os bigodes vão até o
fold mais extremo a no máximo 1,5 vez a distância interquartil da caixa, e os folds além deles não são desenhados.
O losango marca a média, que fica fora da caixa quando há folds atípicos, como na precisão de
\texttt{kstar\_classes}. GB: gaussiano; KNN: $k$-NN; PZ: Parzen; RL: regressão logística; VM: voto
majoritário.}\label{fig:boxplots}
\end{figure}

Na base original, os classificadores formam dois patamares: a regressão logística e o voto majoritário têm erro próximo de 7\,\% 
e F-measure de 0,923, com o $k$-NN logo abaixo (erro de 8,5\,\%), enquanto o gaussiano e a janela de Parzen têm erro próximo 
de 17\,\%. Entre os dois melhores, o voto majoritário tem a maior precisão e a regressão logística a maior cobertura. Na versão \texttt{kstar\_classes}, 
as taxas de erro são todas muito menores que o erro trivial de 0,74\,\% (o melhor modelo erra, em média, cerca de 0,35 exemplo por fold de teste), 
e as diferenças aparecem sobretudo na cobertura: o gaussiano reconhece pior a classe rara (cobertura macro de 0,87), enquanto a regressão logística 
quase não a perde (0,995). Os intervalos são estreitos porque refletem apenas a variação devida à escolha das partições: como as 30 repetições usam a 
mesma base, eles não incorporam a variação que se teria com uma nova amostra de e-mails, e não existe estimador não viesado dessa variância a partir de 
uma única validação cruzada~\cite{bengio2004noUnbiased}. As amplitudes maiores em \texttt{kstar\_classes} decorrem dos 3 ou 4 exemplos da classe rara 
por fold de teste, que tornam as métricas macro muito discretas.

\subsection{Testes de Friedman e Nemenyi}\label{sec:friedman}

Para cada versão da base e cada métrica, os cinco classificadores foram comparados pelo teste de Friedman~\cite{friedman1937ranks}, com os 300 folds 
externos como blocos. Em cada bloco, os classificadores recebem postos de 1 (melhor: menor erro ou maior precisão, cobertura ou F-measure) a 5, 
com postos médios em caso de empate. Com $N=300$ blocos, $k=5$ classificadores e $R_j$ o posto médio do classificador $j$, a estatística
\begin{equation}
\chi^2_F=\frac{12N}{k(k+1)}\Big[\sum_{j=1}^{k}R_j^2-\frac{k(k+1)^2}{4}\Big]
\end{equation}
segue aproximadamente a distribuição $\chi^2$ com $k-1=4$ graus de liberdade sob a hipótese nula de desempenho
igual (usou-se \texttt{scipy.stats.friedmanchisquare}, que corrige os empates). Quando a hipótese nula é rejeitada
com $\alpha=0{,}05$, aplica-se o pós-teste de Nemenyi~\cite{nemenyi1963distribution,demsar2006statistical}:
dois classificadores diferem se seus postos médios diferem mais que a diferença crítica
\begin{equation}
\mathrm{CD}=q_{\alpha}\sqrt{\frac{k(k+1)}{6N}}=\num{2.728}\sqrt{\frac{30}{1800}}\approx\num{0.352},
\end{equation}
com $q_{0{,}05}=\num{2.728}$ para $k=5$~\cite{demsar2006statistical}, e os p-valores par a par foram calculados
com \texttt{scikit-posthocs}.

\begin{table}[H]
\centering
\caption{Teste de Friedman ($N=300$ blocos, 4 graus de liberdade; valor crítico $\chi^2_{0{,}95;4}=\num{9.49}$).}\label{tab:friedman}
\small
\begin{tabular}{@{}l l S[table-format=4.2] c c@{}}
\toprule
Versão & Métrica & {$\chi^2_F$} & p-valor & Decisão ($\alpha=0{,}05$) \\
\midrule
\multirow{4}{*}{\texttt{original\_2\_classes}}
 & Taxa de erro & 1017.59 & $5{,}5\times10^{-219}$ & rejeita $H_0$ \\
 & Precisão & 1033.16 & $2{,}3\times10^{-222}$ & rejeita $H_0$ \\
 & Cobertura & 1047.62 & $1{,}7\times10^{-225}$ & rejeita $H_0$ \\
 & F-measure & 1008.84 & $4{,}3\times10^{-217}$ & rejeita $H_0$ \\
\midrule
\multirow{4}{*}{\texttt{kstar\_classes}}
 & Taxa de erro & 265.55 & $2{,}9\times10^{-56}$ & rejeita $H_0$ \\
 & Precisão & 176.75 & $3{,}7\times10^{-37}$ & rejeita $H_0$ \\
 & Cobertura & 307.65 & $2{,}4\times10^{-65}$ & rejeita $H_0$ \\
 & F-measure & 298.77 & $2{,}0\times10^{-63}$ & rejeita $H_0$ \\
\bottomrule
\end{tabular}
\end{table}

O teste de Friedman rejeitou a igualdade em todas as oito combinações (Tabela~\ref{tab:friedman}), e o pós-teste foi aplicado em todas. 
A Tabela~\ref{tab:postos} traz os postos médios, e a Tabela~\ref{tab:nemenyi} os p-valores de Nemenyi para a F-measure.

\begin{table}[H]
\centering
\caption{Postos médios no teste de Friedman (1 = melhor), com o menor posto de cada linha em negrito; diferenças
maiores que $\mathrm{CD}=\num{0.352}$ são significativas pelo teste de Nemenyi ($\alpha=0{,}05$).}\label{tab:postos}
\small
\begin{tabular*}{\linewidth}{@{\extracolsep{\fill}}l l c c c c c@{}}
\toprule
Versão & Métrica & Gaussiano & $k$-NN & Parzen & Reg.\ logística & Voto maj. \\
\midrule
\multirow{4}{*}{\texttt{original\_2\_classes}}
 & Taxa de erro & \num{4.557} & \num{2.755} & \num{4.443} & \num{1.747} & \textbf{\num{1.498}} \\
 & Precisão & \num{4.297} & \num{2.728} & \num{4.703} & \num{1.972} & \textbf{\num{1.300}} \\
 & Cobertura & \num{4.077} & \num{2.745} & \num{4.923} & \textbf{\num{1.498}} & \num{1.757} \\
 & F-measure & \num{4.387} & \num{2.758} & \num{4.613} & \num{1.702} & \textbf{\num{1.540}} \\
\midrule
\multirow{4}{*}{\texttt{kstar\_classes}}
 & Taxa de erro & \num{3.982} & \num{2.692} & \num{3.010} & \num{2.663} & \textbf{\num{2.653}} \\
 & Precisão & \num{3.817} & \num{2.770} & \num{2.983} & \num{2.842} & \textbf{\num{2.588}} \\
 & Cobertura & \num{4.080} & \num{2.650} & \num{3.027} & \textbf{\num{2.532}} & \num{2.712} \\
 & F-measure & \num{4.063} & \num{2.643} & \num{3.033} & \textbf{\num{2.548}} & \num{2.712} \\
\bottomrule
\end{tabular*}
\end{table}

\begin{table}[H]
\centering
\caption{p-valores do pós-teste de Nemenyi para a F-measure macro, com os pares sem diferença significativa ($p\ge0{,}05$) em negrito. 
GB: gaussiano; KNN: $k$-NN; PZ: Parzen; RL: regressão logística; VM: voto majoritário.}\label{tab:nemenyi}
\small
\setlength{\tabcolsep}{12pt}
\begin{subtable}{\linewidth}
\centering
\caption{\texttt{original\_2\_classes}}
\begin{tabular}{@{}l c c c c c@{}}
\toprule
 & GB & KNN & PZ & RL & VM \\
\midrule
GB & --- & $<\!10^{-4}$ & \ns{0.3998} & $<\!10^{-4}$ & $<\!10^{-4}$ \\
KNN &  & --- & $<\!10^{-4}$ & $<\!10^{-4}$ & $<\!10^{-4}$ \\
PZ &  &  & --- & $<\!10^{-4}$ & $<\!10^{-4}$ \\
RL &  &  &  & --- & \ns{0.7205} \\
VM &  &  &  &  & --- \\
\bottomrule
\end{tabular}
\end{subtable}

\vspace{8pt}
\begin{subtable}{\linewidth}
\centering
\caption{\texttt{kstar\_classes}}
\begin{tabular}{@{}l c c c c c@{}}
\toprule
 & GB & KNN & PZ & RL & VM \\
\midrule
GB & --- & $<\!10^{-4}$ & $<\!10^{-4}$ & $<\!10^{-4}$ & $<\!10^{-4}$ \\
KNN &  & --- & \num{0.0212} & \ns{0.9481} & \ns{0.9844} \\
PZ &  &  & --- & \num{0.0016} & \ns{0.0925} \\
RL &  &  &  & --- & \ns{0.7127} \\
VM &  &  &  &  & --- \\
\bottomrule
\end{tabular}
\end{subtable}
\end{table}

Para as demais métricas, os pares sem diferença significativa foram:
\begin{itemize}
\item \texttt{original\_2\_classes}: taxa de erro, GB--PZ ($p=\num{0.905}$) e RL--VM ($p=\num{0.305}$); precisão,
  nenhum (o voto majoritário é significativamente mais preciso que a regressão logística, $p<10^{-5}$); cobertura,
  RL--VM ($p=\num{0.265}$). Assim, gaussiano e Parzen empatam na taxa de erro e na F-measure, mas o gaussiano tem
  precisão ($p=\num{0.014}$) e cobertura ($p<10^{-9}$) macro significativamente maiores.
\item \texttt{kstar\_classes}: taxa de erro e precisão, KNN--PZ, KNN--RL, KNN--VM, PZ--RL e RL--VM; cobertura,
  KNN--RL, KNN--VM, PZ--VM e RL--VM. Em todas as métricas, o gaussiano difere de todos os demais.
\end{itemize}

Em resumo, na base original há três grupos estatisticamente distintos pela F-measure: \{voto majoritário,
regressão logística\} no topo, o $k$-NN no meio e \{gaussiano, Parzen\} abaixo. Na versão \texttt{kstar\_classes},
regressão logística, $k$-NN e voto majoritário não se distinguem, a janela de Parzen fica abaixo da regressão
logística e do $k$-NN, e o gaussiano é o pior em todas as métricas.

Na interpretação desses resultados, é preciso considerar, em primeiro lugar, que o roteiro de
Demšar~\cite{demsar2006statistical} foi proposto para comparar classificadores em várias bases independentes,
enquanto aqui os blocos são folds de uma única base, correlacionados entre si, e o grande número de blocos ($N=300$)
torna o teste muito sensível. Os p-valores extremamente pequenos indicam, portanto, que a ordenação dos
classificadores é muito estável entre as partições da Spambase, e não que ela se generalize a outras bases. Em
segundo lugar, comparações por postos médios dependem do conjunto de classificadores
incluído~\cite{benavoli2016posthoc}, e por isso as conclusões acima são sempre lidas junto com as diferenças nas
próprias métricas (Tabelas~\ref{tab:ic-original} e~\ref{tab:ic-kstar}).

\subsection{Curvas de aprendizagem}\label{sec:curvas}

\paragraph{Protocolo.} Para cada versão da base e cada proporção de treino $p\in\{5\%,10\%,\dots,95\%\}$ (19
valores), a base completa foi dividida 10 vezes em treino ($p$) e teste ($1-p$) por amostragem estratificada
(\texttt{StratifiedShuffleSplit}, semente 42). Em cada divisão, os hiperparâmetros dos quatro classificadores
de base foram escolhidos por validação cruzada estratificada usando apenas a parte de treino, com as mesmas grades
e o mesmo critério da Seção~\ref{sec:hiper}; o voto majoritário usa os membros ajustados. O número de folds
internos é $\min(5,n_{\min})$, em que $n_{\min}$ é o tamanho da menor classe no treino: na versão
\texttt{kstar\_classes}, com 5\,\% e 10\,\% de treino há só 2 e 3 exemplos da classe rara, e a CV interna usa 2 e 3
folds. Calculou-se a F-measure macro tanto no conjunto de treino (resubstituição) quanto no de teste, totalizando
$2\times19\times10\times5=\num{1900}$ avaliações. As Figuras~\ref{fig:lc-original} e~\ref{fig:lc-kstar} mostram
média $\pm$ 1 desvio-padrão das 10 divisões, e a Tabela~\ref{tab:lc} alguns pontos das curvas.

\begin{table}[H]
\centering
\caption{F-measure macro média de treino / teste em algumas proporções de treino (10 divisões por proporção).}\label{tab:lc}
\footnotesize
\setlength{\tabcolsep}{4pt}
\begin{tabular*}{\linewidth}{@{\extracolsep{\fill}}l c c c c c c@{}}
\toprule
Classificador & 5\,\% & 10\,\% & 25\,\% & 50\,\% & 75\,\% & 95\,\% \\
\midrule
\multicolumn{7}{@{}l}{\texttt{original\_2\_classes} ($n_{\text{treino}}$: 230, 460, \num{1150}, \num{2300}, \num{3450}, \num{4370})} \\
Bayesiano gaussiano & \lc{0.904}{0.842} & \lc{0.872}{0.836} & \lc{0.843}{0.830} & \lc{0.833}{0.829} & \lc{0.830}{0.829} & \lc{0.831}{0.829} \\
Bayesiano $k$-NN & \lc{0.932}{0.831} & \lc{0.929}{0.854} & \lc{0.941}{0.877} & \lc{0.961}{0.896} & \lc{0.942}{0.909} & \lc{0.943}{0.898} \\
Janela de Parzen & \lc{0.915}{0.683} & \lc{0.930}{0.710} & \lc{0.991}{0.766} & \lc{0.985}{0.795} & \lc{0.986}{0.810} & \lc{0.993}{0.814} \\
Regressão logística & \lc{0.950}{0.883} & \lc{0.936}{0.898} & \lc{0.928}{0.912} & \lc{0.928}{0.918} & \lc{0.926}{0.919} & \lc{0.929}{0.904} \\
Voto majoritário & \lc{0.966}{0.871} & \lc{0.959}{0.888} & \lc{0.964}{0.904} & \lc{0.968}{0.916} & \lc{0.958}{0.920} & \lc{0.959}{0.910} \\
\midrule
\multicolumn{7}{@{}l}{\texttt{kstar\_classes} (exemplos da classe rara no treino: 2, 3, 8, 17, 25, 32)} \\
Bayesiano gaussiano & \lc{1.000}{0.569} & \lc{1.000}{0.561} & \lc{1.000}{0.755} & \lc{0.994}{0.873} & \lc{0.985}{0.893} & \lc{0.974}{0.873} \\
Bayesiano $k$-NN & \lc{0.950}{0.895} & \lc{1.000}{0.956} & \lc{0.997}{0.950} & \lc{0.995}{0.971} & \lc{0.997}{0.968} & \lc{1.000}{0.990} \\
Janela de Parzen & \lc{0.966}{0.833} & \lc{1.000}{0.878} & \lc{1.000}{0.904} & \lc{1.000}{0.930} & \lc{1.000}{0.947} & \lc{1.000}{0.958} \\
Regressão logística & \lc{0.866}{0.869} & \lc{1.000}{0.927} & \lc{0.997}{0.950} & \lc{1.000}{0.964} & \lc{1.000}{0.962} & \lc{1.000}{0.980} \\
Voto majoritário & \lc{0.950}{0.789} & \lc{1.000}{0.888} & \lc{1.000}{0.914} & \lc{1.000}{0.946} & \lc{1.000}{0.972} & \lc{1.000}{0.966} \\
\bottomrule
\end{tabular*}
\end{table}

\begin{figure}[!tbp]
\centering
\lcfigure{original_2_classes}{0.6}{0.6,0.7,...,1.0}
\caption{Curvas de aprendizagem na versão \texttt{original\_2\_classes}: F-measure macro de treino e de teste,
média das 10 divisões estratificadas de cada proporção de treino, com faixa de $\pm1$ desvio-padrão (limitada a 1)
e a mesma escala em todos os painéis.}\label{fig:lc-original}
\end{figure}

\begin{figure}[!tbp]
\centering
\lcfigure{kstar_classes}{0.4}{0.4,0.5,...,1.0}
\caption{Curvas de aprendizagem na versão \texttt{kstar\_classes}, com as mesmas convenções da
Figura~\ref{fig:lc-original} e escala vertical a partir de 0,4.}\label{fig:lc-kstar}
\end{figure}

\paragraph{Comentários: \texttt{original\_2\_classes}.} As curvas separam os classificadores pelo compromisso entre
viés e variância~\cite{hastie2009elements}:
\begin{itemize}
\item \textbf{Bayesiano gaussiano (viés alto).} As curvas de treino e teste convergem já com cerca de 25\,\% dos
  dados para o mesmo patamar, $F\approx0{,}83$. A F-measure de teste é, na verdade, máxima com 5\,\% de treino
  (0,842) e cai levemente até 0,83, sem melhorar depois disso. Mais dados não ajudam, porque o erro vem do modelo: a
  normal multivariada descreve mal atributos esparsos, com excesso de zeros e caudas longas.
\item \textbf{Janela de Parzen (variância alta).} A F-measure de treino fica perto de 0,99, enquanto a de teste
  parte de 0,68 e chega a 0,81. A grande distância entre as curvas indica sobreajuste: em 57 dimensões, um
  estimador de densidade não paramétrico precisa de muitos exemplos~\cite{silverman1986density}. É também o
  modelo que mais ganha com o aumento da amostra (13 pontos percentuais entre 5\,\% e 90\,\%), e a curva de teste
  ainda sobe no fim.
\item \textbf{$k$-NN.} A curva de teste cresce de 0,83 para cerca de 0,90--0,91, com distância entre treino e teste
  que diminui aos poucos: o método se beneficia do adensamento da amostra.
\item \textbf{Regressão logística e voto majoritário.} São os melhores desde as menores amostras ($F$ de teste de
  0,88 e 0,87 com 5\,\%) e estabilizam perto de 0,92 a partir de cerca de 50\,\%. A regressão logística tem a menor
  distância entre treino e teste e o padrão clássico de um modelo linear regularizado: a F-measure de treino cai de
  0,950 para 0,929 enquanto a de teste sobe. Em 95\,\%, a F-measure de teste cai um pouco e varia mais, porque o
  teste tem apenas 231 exemplos.
\end{itemize}

\paragraph{Comentários: \texttt{kstar\_classes}.} A tarefa de separar os 34 e-mails do grupo $P_2$ é mais fácil
para os modelos locais e lineares: $k$-NN e regressão logística superam $F=0{,}92$ no teste com apenas 10\,\% dos
dados (3 exemplos da classe rara) e chegam a 0,98--0,99 no fim. O gaussiano é o que mais depende do tamanho da
amostra: com 2 ou 3 exemplos da classe rara, a covariância estimada tem posto no máximo 1 ou 2 e a regularização
$\varepsilon\mathbf{I}$ domina as demais direções; o modelo acerta todo o treino ($F=1$) e generaliza mal
($F\approx0{,}56$). A F-measure de teste sobe rapidamente até cerca de 50\,\% e depois oscila entre 0,83 e
0,93, com a maior distância para o treino entre os modelos, enquanto a janela de Parzen melhora, com oscilações, de
0,83 para 0,96. O voto majoritário parte de 0,79 com 5\,\%, abaixo de três dos seus membros, provavelmente porque, com
o gaussiano errando muito a classe rara, os empates de 2 a 2 são resolvidos a favor da classe 0; ele alcança 0,97
com 75\,\%. Nas maiores proporções, o teste contém só 2 a 9 exemplos da classe rara, o que explica as faixas largas de
desvio-padrão.

%======================================================================
\section{Discussão e limitações}\label{sec:discussao}
%======================================================================

A primeira limitação diz respeito à escolha de $K^*$, pois a silhueta selecionou uma partição dominada por 34 objetos atípicos
(Seção~\ref{sec:q1-discussao}), sem relação com as classes a priori, e é provável que uma transformação que reduza as caudas dos atributos 
antes do agrupamento, como $\log(1+x)$, ou um índice de validação menos sensível a grupos pequenos e isolados levasse a outra escolha; manteve-se, 
ainda assim, o critério definido no enunciado. Essa escolha determina a natureza da versão \texttt{kstar\_classes}: como os rótulos $y^*$ foram 
obtidos pelo agrupamento da base completa, a tarefa de  classificação consiste em reproduzir uma regra de agrupamento conhecida, e não em prever 
um fenômeno externo. Essa regra é um hiperplano no espaço padronizado (Equação~\ref{eq:fronteira}), o que favorece a regressão logística e 
o $k$-NN euclidiano, e, como a classe rara tem apenas 34 exemplos, as métricas macro de cada fold dependem de 3 ou 4 exemplos.

Algumas decisões de implementação dos classificadores também afetam os resultados, a começar pela janela de Parzen, cujo estimador recebe os atributos 
na escala original, com a mesma janela $h$ em todas as dimensões. Como os desvios-padrão variam de $\num{0.076}$ a $\num{606}$, a distância no kernel 
é dominada pelos atributos \texttt{capital\_run\_length\_*}, e uma janela de 1 ou 2 é larga demais para as frequências e estreita demais para as
contagens de maiúsculas, o que provavelmente contribui para o desempenho fraco e o sobreajuste observados; a correção natural seria padronizar os 
atributos dentro do modelo, como no $k$-NN. No classificador gaussiano, o termo $\varepsilon\mathbf{I}$ é um desvio mínimo, mas real, da estimativa 
de máxima verossimilhança pedida no enunciado, que se justifica porque a densidade não está definida quando $\hat{\bm\Sigma}_i$ é singular. 
Na regressão logística, $C=100$, o maior valor da grade, foi escolhido em 160 dos 300 folds da base original, de modo que uma grade mais ampla 
permitiria verificar se uma regularização ainda mais fraca traria ganho. Por fim, no voto majoritário, os empates de 2 a 2 entre os quatro
votantes vão para a classe 0; na base original, essa regra favorece \enquote{não spam} e é compatível com o comitê ter a maior precisão e uma 
cobertura menor que a da regressão logística.

No protocolo de avaliação, por fim, as 391 linhas duplicadas foram mantidas, e uma cópia no treino e outra no teste é acertada trivialmente pelo $k$-NN 
com $k=1$, escolhido em 73 folds da base original e em 242 da versão \texttt{kstar\_classes}; essas estimativas são, portanto, levemente otimistas, 
sobretudo na classe rara de \texttt{kstar\_classes}, em que só 28 dos 34 e-mails são distintos. Além disso, tanto os intervalos de confiança
quanto os testes de Friedman e Nemenyi usam resultados de partições da mesma base e, por isso, subestimam a incerteza em relação a novas amostras 
(Seções~\ref{sec:ic} e~\ref{sec:friedman}), algo que alternativas como o teste $t$ corrigido de Nadeau e Bengio~\cite{nadeau2003inference} levam
em conta; ao avaliarem a replicabilidade de vários testes, Bouckaert e Frank~\cite{bouckaert2004evaluating} optaram justamente por ele. Finalmente, seguindo o item d) do enunciado, as curvas de aprendizagem foram calculadas apenas para a F-measure, e as demais métricas não 
foram registradas nesse experimento.

%======================================================================
\section{Conclusões}\label{sec:conclusao}
%======================================================================

Na Questão~1, as 900 execuções do double k-means convergiram pelo critério dos slides, e a silhueta das melhores
execuções de cada par selecionou $(K^*,H^*)=(2,2)$, com $\Sil=\num{0.6596}$, mais que o dobro do segundo maior
valor. A partição obtida isola 34 e-mails legítimos que usam com frequência muito alta 11 palavras e números ligados
à HP, origem dos e-mails não spam da base, e tem $\ARI=\num{-0.00495}$ frente às classes spam/não spam, o que mostra
que a estrutura geométrica dominante na base padronizada não é a divisão entre spam e não spam, mas a separação de
um pequeno grupo atípico. Esse resultado decorre tanto do critério de seleção quanto do próprio algoritmo, já que a
silhueta, por ser uma média sobre todos os objetos, premia a separação de poucos pontos extremos, e a soma de
quadrados do double k-means dá peso desproporcional às caudas longas dos atributos esparsos. Além disso, a análise
da partição na convergência mostrou que a fronteira entre os dois grupos é um hiperplano no espaço padronizado
(Equação~\ref{eq:fronteira}), fato que antecipa o comportamento dos classificadores na segunda versão da base.

Na Questão~2, a validação cruzada estratificada $30\times10$ com ajuste interno de hiperparâmetros, aplicada da
mesma forma aos cinco classificadores, mostrou que, na detecção de spam, o voto majoritário ($F=\num{0.9238}$,
erro de 7,2\,\%) e a regressão logística ($F=\num{0.9230}$, erro de 7,3\,\%) são estatisticamente equivalentes e
superiores aos demais, seguidos pelo $k$-NN com distância city-block ($F=\num{0.9093}$), enquanto o gaussiano e a
janela de Parzen ficaram abaixo ($F\approx0{,}82$--$0{,}83$). Embora empatem na F-measure, os dois melhores modelos
chegam a esse resultado por caminhos diferentes: o voto majoritário tem a maior precisão e a regressão logística a
maior cobertura, diferença coerente com a regra de desempate do comitê a favor da classe não spam. As curvas de
aprendizagem ajudam a explicar a ordenação, pois indicam que o gaussiano é limitado por viés, já que a normal
multivariada descreve mal atributos esparsos e de cauda longa e mais dados não melhoram o seu desempenho, ao passo
que a janela de Parzen é limitada por variância, com grande distância entre treino e teste em 57 dimensões e uma
curva de teste que ainda sobe com 95\,\% dos dados. A regressão logística, por sua vez, tem bom desempenho desde as
menores amostras e a menor distância entre treino e teste, o que, diante do empate estatístico com o voto
majoritário, que exige ajustar quatro modelos, a torna a opção mais econômica para essa tarefa.

Na versão com os rótulos da Questão~1, regressão logística, $k$-NN e voto majoritário reproduzem quase perfeitamente
a partição ($F$ entre $\num{0.967}$ e $\num{0.976}$) e não se distinguem estatisticamente, resultado esperado porque a
fronteira entre os grupos é linear e foi definida pela distância euclidiana nos atributos padronizados, a mesma
geometria do $k$-NN escolhido na maioria dos folds. O gaussiano é o pior em todas as métricas, prejudicado por uma
classe rara com menos exemplos de treino que atributos, cuja covariância estimada é singular sem o termo
$\varepsilon\mathbf{I}$. Como essa versão mede a capacidade de reproduzir uma regra de agrupamento, e não de prever
um fenômeno externo, os seus resultados não devem ser lidos como evidência da utilidade prática dos classificadores.

Do ponto de vista metodológico, a separação entre seleção e avaliação mostrou-se necessária, pois, na versão
\texttt{kstar\_classes}, a F-measure interna da combinação escolhida para o $k$-NN ($\num{0.9879}$) supera a de teste
($\num{0.9722}$), um otimismo que teria sido reportado como desempenho caso a mesma validação servisse para escolher
e para avaliar os modelos. A comparação pareada nas mesmas 300 partições, por sua vez, deu aos testes de Friedman e
Nemenyi grande poder para ordenar os classificadores, embora os p-valores reflitam a estabilidade dessa ordenação
dentro da Spambase, e não a sua generalização para outras bases. Como continuação, as limitações da
Seção~\ref{sec:discussao} apontam caminhos diretos: repetir o agrupamento após uma transformação que reduza as caudas
dos atributos, padronizar os atributos dentro da janela de Parzen, ampliar a grade de $C$ da regressão logística e
substituir os intervalos e testes por procedimentos que considerem a dependência entre folds, como o teste $t$
corrigido de Nadeau e Bengio~\cite{nadeau2003inference}.

% Printed in English so biblatex uses the IEEE strings (vol., no., pp., and, Eds., [Online]. Available).
\begin{otherlanguage}{american}
\printbibliography[heading=bibintoc,title={Referências}]
\end{otherlanguage}

\appendix
%======================================================================
\section{Reprodução dos resultados}\label{app:reproducao}
%======================================================================

Os experimentos usaram o arquivo \texttt{dataset/spambase.data} do repositório UCI, cujo resumo SHA-256 é
\begin{center}
\small\texttt{b1ef93de71f97714d3d7d4f58fc9f718da7bbc8ac8a150eff2778616a8097b12}.
\end{center}
Na raiz do repositório, com as dependências instaladas por \texttt{poetry install}, os comandos abaixo, nesta
ordem, geram todos os resultados usados no relatório:

\begin{enumerate}[leftmargin=*]
\item \texttt{poetry run python -m src.experiments.question\_1}\\
  900 execuções do double k-means e melhor execução de cada par $(K,H)$ (\texttt{results/q1\_*}).
\item \texttt{poetry run python -m src.analysis.silhouette}\\
  Silhueta dos 9 vencedores e escolha de $(K^*,H^*)$ (\texttt{results/silhouette\_scores.csv}).
\item \texttt{poetry run python -m src.analysis.clustering\_evaluation}\\
  ARI, matriz $\mathbf{G}$, contingência e rótulos $y^*$ (\texttt{results/q1\_object\_labels.csv}).
\item \texttt{poetry run python -m src.experiments.question\_2 -{}-n-jobs -1}\\
  Validação cruzada $30\times10$ com ajuste interno (\texttt{results/q2/cv\_30x10\_results.csv}).
\item \texttt{poetry run python -m src.analysis.statistical\_analysis}\\
  IC 95\,\%, Friedman e Nemenyi (arquivos em \texttt{results/}).
\item \texttt{poetry run python docs/relatorio/gerar\_boxplots.py}\\
  Estatísticas dos \textit{boxplots} da Figura~\ref{fig:boxplots} (\texttt{docs/relatorio/dados/boxplots.csv}).
\item \texttt{poetry run python -m src.experiments.learning\_curves -{}-n-jobs -1}\\
  Curvas de aprendizagem (arquivos em \texttt{results/} e \texttt{plots/}).
\end{enumerate}

Para compilar este relatório: \texttt{latexmk -pdf relatorio.tex} no diretório \texttt{docs/relatorio}.

\end{document}
